\section{Design Goals and Non-Goals}
\label{sec:goals}

\subsection{Goals}

\begin{description}[leftmargin=1.2em]
  \item[G1. No custody.] No party other than the buyer and the vendor ever holds a key, or a set of keys, sufficient to move a buyer's payment. The moderator in particular holds no signing key over escrowed funds.
  \item[G2. Verifiable public catalog.] Anyone can check that a listing was published by the store it names and has not been altered. A listing's address is derived from its content.
  \item[G3. Private orders.] An order's contents, including any ship-to address, are readable only by the vendor it is addressed to. No public relay or index stores a record that a given buyer wrote to a given vendor.
  \item[G4. Bounded moderation.] A moderator can choose between two payouts that buyer and vendor already signed. It cannot sign alone, redirect funds, change amounts, or learn about an order unless a party opens a dispute.
  \item[G5. Liveness without a platform.] A buyer can recover escrowed funds from a vendor who disappears, and a vendor can collect from a buyer who stops responding, without anyone's cooperation beyond broadcasting a presigned transaction.
  \item[G6. Thin clients.] A phone client embeds Tor so it can reach the node it is paired with. It does not hold spend keys, and it does not dial other vendors. A browser on the same machine talks to the node on a loopback address. Spending and chain watches stay on the node.
  \item[G7. Honest degradation.] When a privacy mechanism is unavailable (no Tor, no Payjoin partner, no Lightning route), the system says so rather than silently falling back.
\end{description}

\subsection{Non-goals}

\begin{itemize}
  \item \textbf{Anonymity against a global passive adversary.} We do not claim unlinkability against an adversary that observes most of the network. Section~\ref{sec:privacy} states what such an adversary learns.
  \item \textbf{Hiding the catalog.} Listings are public by design. A market that cannot be discovered does not serve its purpose.
  \item \textbf{Hiding the ship-to address from the seller or the carrier.} A seller of a physical good must learn where to send it, and the carrier must deliver it. This is a ceiling on buyer privacy for physical goods, not a gap to be engineered away.
  \item \textbf{Deciding factual disputes cryptographically.} Whether a parcel contained the right item is a judgment by a person. The protocol bounds what that person can do with the funds; it does not replace the judgment.
  \item \textbf{Reputation.} A review or rating system is outside this paper.
  \item \textbf{Content policy.} The protocol is designed for ordinary goods. Policies for what may be listed, and the obligations of people who operate replicas or seeds, are outside the scope of this paper.
\end{itemize}
